% Aesthetic Wing concept sheet — speculative, non-confidential.
% Build from this directory: xelatex -interaction=nonstopmode sheet.tex
\documentclass[10pt]{article}
\usepackage[margin=0.62in]{geometry}
\usepackage[table]{xcolor}
\usepackage{graphicx}
\usepackage{tabularx}
\usepackage{array}
\usepackage{booktabs}
\usepackage{fontspec}
\usepackage{tikz}
\usepackage{eso-pic}
\pagestyle{empty}
\setlength{\parindent}{0pt}
\setlength{\parskip}{0.42em}
\renewcommand{\arraystretch}{1.14}

\definecolor{ink}{HTML}{202126}
\definecolor{paper}{HTML}{FBF7EB}
\definecolor{line}{HTML}{DED6C2}
\definecolor{wingmint}{HTML}{493A72}
\definecolor{wingcoral}{HTML}{B63A83}
\definecolor{wingblue}{HTML}{3955C6}
\definecolor{wingsun}{HTML}{7CA634}

\newfontfamily\processing[
  Path=../../../../system/public/type/webfonts/,
  Extension=.ttf,
  UprightFont=ywft-processing-regular,
  BoldFont=ywft-processing-bold
]{YWFT Processing}
\newfontfamily\berkeley[
  Path=../../../../bills/invoices/,
  Extension=.ttf,
  UprightFont=BerkeleyMonoVariable-Regular,
  ItalicFont=BerkeleyMonoVariable-Italic
]{Berkeley Mono}
\setmainfont{Avenir Next}

\newcommand{\sheetid}{AC-AWG-01}
\newcommand{\docdate}{Concept sheet · July 23, 2026}
\newcommand{\hero}{../gens/v6-more-lift.png}
\newcommand{\roll}{../gens/v8-road-mode.png}
\newcommand{\step}{../gens/v5-wing-wakes.png}
\newcommand{\pals}{../../../../bills/invoices/pals.pdf}

\AddToShipoutPictureBG{%
  \begin{tikzpicture}[remember picture,overlay]
    \fill[paper] (current page.south west) rectangle (current page.north east);
    \draw[line,line width=0.45pt]
      ([xshift=0.62in,yshift=0.55in]current page.south west) --
      ([xshift=-0.62in,yshift=0.55in]current page.south east);
    \node[anchor=south,inner sep=0pt] at ([yshift=0.35in]current page.south)
      {\makebox[\dimexpr\paperwidth-1.24in\relax]{\berkeley\scriptsize\color{ink}
      \docdate\hfill\sheetid\hfill\thepage\,/\,2}};
  \end{tikzpicture}%
}

\newcommand{\titleblock}[2]{%
  \begin{tabularx}{\linewidth}{@{}Xr@{}}
    {\processing\fontsize{27pt}{28pt}\bfseries\color{#1} Aesthetic Wing} &
    \raisebox{-0.45cm}{
\includegraphics[height=1.25cm]{\pals}}\\
    {\processing\fontsize{16pt}{18pt}\bfseries\color{ink} #2} & {}
  \end{tabularx}
  \vspace{0.35em}\color{#1}\rule{\linewidth}{6pt}\color{ink}\vspace{0.55em}
}

\newcommand{\head}[2]{%
  \vspace{0.38em}{\processing\fontsize{14pt}{16pt}\bfseries\color{#1}#2}\par
}

\newcommand{\tag}[1]{%
  \tikz[baseline=-0.6ex]{\node[fill=wingmint!12,draw=wingmint!55,
    rounded corners=4pt,inner xsep=6pt,inner ysep=3pt]{\berkeley\scriptsize #1};}}

\begin{document}
\color{ink}

\titleblock{wingmint}{Real chill street aviation}

\begin{tabularx}{\linewidth}{@{}Xr@{}}
  \textbf{Model:} \sheetid & \textbf{Status:} speculative research concept\\
  \textbf{Class:} personal electric mobility / bounded lift &
  \textbf{Phase I:} unmanned only
\end{tabularx}

\vspace{0.45em}
\begin{minipage}[t]{0.54\linewidth}
\vspace{0pt}
\head{wingmint}{The short version}
Aesthetic Wing is an electric one-wheel that can make an \textbf{air step}.
Most of the time it rolls. When the surface changes, two guarded lift petals
open, progressively unload the wheel, rise through a bounded envelope, and
settle back onto the street.

It is not a flying car. It proposes a smaller gesture between riding and flying:
\textbf{a little more street, in every direction.}

\head{wingcoral}{The research question}
Can a compact one-wheel vehicle transition safely between rolling, partial lift,
and controlled low hover while preserving an intuitive board-riding experience?

\head{wingblue}{The interface}
\tag{lean to move}\quad\tag{press to rise}\quad\tag{relax to settle}

The rider expresses intent through stance and foot pressure. Stabilization,
height limiting, energy reserve, and automatic settling remain machine duties.

\head{wingsun}{Truth before spectacle}
The illustration is a product vision, not flight-test evidence. Phase I contains
no rider, no public-street operation, and no unrestricted altitude.
\end{minipage}\hfill
\begin{minipage}[t]{0.42\linewidth}
\vspace{0pt}
\begin{tikzpicture}
  \node[inner sep=0pt] (im) {\includegraphics[width=\linewidth]{\hero}};
  \draw[rounded corners=9pt,line,line width=0.8pt]
    (im.south west) rectangle (im.north east);
\end{tikzpicture}

{\berkeley\scriptsize\color{ink!70}Concept illustration: Jeffrey in a controlled eight-foot hover. The vehicle depicted does not exist.}
\end{minipage}

\vfill
\newpage

\titleblock{wingblue}{Phase I research program}

\head{wingblue}{Operating modes}
\begin{tabularx}{\linewidth}{>{\bfseries}p{1.1in}X}
Roll & One central wheel supports ordinary efficient travel.\\
Skim & Lift partially unloads the wheel over rough ground.\\
Air step & A short controlled rise, translation, and landing.\\
Hover & Fully airborne movement inside a deliberately low envelope.\\
Safe settle & Automatic descent after lost input, low reserve, or fault.\\
\end{tabularx}

\vspace{0.4em}
\begin{minipage}[t]{0.49\linewidth}
\vspace{0pt}
\head{wingmint}{Phase I objectives}
\begin{enumerate}
  \setlength{\itemsep}{0.22em}
  \item Model mass, lift area, power, thermal behavior, downwash, noise, and transition dynamics.
  \item Validate guarded propulsion on a stationary full-load test article.
  \item Develop roll-to-hover controls and automatic settling in simulation and hardware.
  \item Demonstrate an unmanned weighted prototype rolling, unloading, crossing a controlled interruption, landing, and rolling again.
\end{enumerate}

\head{wingcoral}{Safety gates}
Rotor containment; independent shutdown; tethered initial trials; battery
isolation; thermal monitoring; motor-out testing; bounded height and speed;
formal hazard review before every expansion of the envelope.
\end{minipage}\hfill
\begin{minipage}[t]{0.47\linewidth}
\vspace{0pt}
\head{wingsun}{Twelve months / \$300,000}
\begin{tabularx}{\linewidth}{@{}p{0.65in}X@{}}
1--2 & requirements, prior art, hazards\\
2--4 & energy and geometry trade study\\
3--6 & guarded-lift test rig\\
5--8 & rolling mule and controller\\
7--10 & integrated tethered tests\\
9--11 & unmanned air-step demonstration\\
11--12 & validation and Phase II design
\end{tabularx}

\vspace{0.55em}
\begin{tikzpicture}
  \node[inner sep=0pt] (a) {\includegraphics[width=0.48\linewidth]{\roll}};
  \node[inner sep=0pt,anchor=west] (b) at ([xshift=5pt]a.east)
    {\includegraphics[width=0.48\linewidth]{\step}};
  \draw[line,line width=0.6pt] (a.south west) rectangle (a.north east);
  \draw[line,line width=0.6pt] (b.south west) rectangle (b.north east);
\end{tikzpicture}

{\berkeley\fontsize{7.4pt}{9pt}\selectfont\color{ink!70}%
\parbox{\linewidth}{Roll → air step. Generated concept sequence.\\
Performance remains to be measured.}}
\end{minipage}

\head{wingmint}{Funding request}
\textbf{\$300,000 over 12 months} to establish technical feasibility and
demonstrate an unmanned wheel-to-hover transition. The non-confidential funding
thesis is a deployable distributed-lift and transition-control system for
bounded personal mobility. The product and cultural identity remain
\textbf{Aesthetic Wing}.

\vfill
{\berkeley\small\color{ink!70}aesthetic.computer · speculative hardware research · not offered for sale}

\end{document}
